% ===== tr-03 块1：第 3 章 pp.22--25（原文页 22--25）=====
% [原文 p.22]
\begin{exercise}{3.6.1}
假设存在一点 $x\in\mathcal X_0$，它的一个邻域同构于某个仿射 Gorenstein 环面簇 $Y_x$ 的一个零维环面轨道的邻域。这样的仿射簇可如下具体给出。置 $M=\Z^n$，$M_\R=M\otimes_\Z\R$，$N=\operatorname{Hom}_\Z(M,\Z)$，$N_\R=N\otimes_\Z\R$，如第 4 节。于是存在格多面体 $\sigma\subseteq M_\R$（其中 $n=\dim\mathcal X_t$），$C(\sigma):=\{(rm,r)\mid m\in\sigma,\ r\ge0\}\subseteq M_\R\oplus\R$，$P:=C(\sigma)^\vee\cap(N\oplus\Z)$ 为由锥 $C(\sigma)$ 的对偶所确定的幺半群，最后 $Y_x=\operatorname{Spec}\C[P]$，而 $f$ 与单项式 $z^{(0,1)}$ 相合。
% [Tutor Note: 原文习题正文为斜体，末端的 □ 是作者标注的习题结束符；$\tilde U_\delta$ 中的 $z^p$ 是幺半群元素 $p\in P$ 对应的单项式（特征标），不是 $z$ 乘 $p$，文字层常抽成 |zp|]

现在取 $x$ 的一个如下形状的小邻域
\[
\tilde U_\delta=\{y\in\operatorname{Spec}\C[P]\mid |z^p|<\delta\ \text{对所有}\ p\in P\}.
\]
这是一个开集，因为条件 $|z^p|<\delta$ 只需在 $P$ 的一个有限生成集上检验，只要 $\delta<1$。然后证明：对给定的 $t$（$|t|<1$）以及 $\epsilon=-2\pi/\log|t|$，若
\[
\sigma_t:=\{m\in M_\R\mid \langle p,(m,1)\rangle>\frac{\log\delta}{\log|t|}\ \text{对所有}\ p\in P\},
\]
则
\[
f^{-1}(t)\cap\tilde U_\delta\cong X_\epsilon(\sigma_t).
\]
注意
\[
\sigma:=\{m\in M_\R\mid \langle p,(m,1)\rangle\ge0\ \text{对所有}\ p\in P\},
\]
于是 $\sigma_t$ 是 $\sigma$ 的开子集，且当 $t\to0$ 时 $\sigma_t$ 收敛到 $\sigma$ 的内部。\proofend
\end{exercise}

这一观察有望为下一节的基本构造提供动机。

\section{环面退化、相交复形及其对偶}

% [Tutor Note: 本节论证顺序是"倒推"------原文明确说 from the opposite direction：先用上一节（§6）的动机引出对偶相交复形，再回溯到对数结构，读者不要误以为对数结构是出发点]
现在回到严格的陈述。我想解释文献 [22] 所启动的纲领背后的基本想法。尽管我会把前面几节用作动机，这项工作实际上发端于 Siebert 开始的一个纲领：通过 Calabi--Yau 流形的退化来研究镜像对称。Schröer 与 Siebert 的工作 [54]、[55] 使 Siebert 意识到，Calabi--Yau 流形退化上的对数结构能让人把镜像对称看作施加于退化 Calabi--Yau 簇之上的一个操作。Siebert 观察到，在组合层面上，镜像对称交换的是关于对数结构的数据与一个极化。这一点将在下一节引入对数结构时解释得更清楚。Siebert 和我一同认识到，他所考虑的组合数据可以自然地编码在退化的对偶相交复形中，而镜像对称于是对应于对偶相交复形上的离散 Legendre 变换。此后就清楚了：这一途径给出了 SYZ 猜想的代数几何化。

这里我将从相反的方向来解释这个纲领：先从上一节引入对偶相交复形的动机出发，然后倒推回去，直到自然到达对数结构。本节的大部分材料来自 [22] 的第 4 节。

\begin{definition}{3.7.1}
设 $f:\mathcal X\to D$ 是一个相对维数为 $n$ 的固有平坦族，其中 $D$ 是圆盘，$\mathcal X$ 是复解析空间（不必非奇异）。我们称 $f$ 是 Calabi--Yau 簇的环面退化，如果
% [原文 p.23]
\begin{enumerate}[label=(\arabic*)]
\item $\mathcal X_t$ 是不可约正规的 Calabi--Yau 簇，且对 $t\ne0$ 只具有典范奇点。（读者不妨假设 $\mathcal X_t$ 对 $t\ne0$ 是光滑的。）
\item 若 $\nu:\tilde{\mathcal X}_0\to\mathcal X_0$ 是正规化，则 $\tilde{\mathcal X}_0$ 是环面簇的不交并，导子轨迹 $C\subseteq\tilde{\mathcal X}_0$ 是既约的，且映射 $C\to\nu(C)$ 非分歧、一般地二对一。（导子轨迹是 $\nu$ 不为同构的集合上自然定义的概形结构。）方块
% [Tutor Note: 下图的左竖直箭头原文无标签，它是 $C$ 到 $\nu(C)$ 的对应映射；右竖直箭头的 $\nu$ 标注在箭头右侧。原文称该方块 cartesian and cocartesian，即它既是纤维积方块又是纤维余积方块]
\[
\begin{tikzcd}
C \arrow[r] \arrow[d] & \tilde{\mathcal X}_0 \arrow[d, "\nu"]\\
\nu(C) \arrow[r] & \mathcal X_0
\end{tikzcd}
\]
是笛卡尔的（即纤维积方块），也是余笛卡尔的（即纤维余积方块）。
\item $\mathcal X_0$ 是既约的 Gorenstein 空间，且导子轨迹 $C$ 限制在 $\tilde{\mathcal X}_0$ 的每个不可约分支上，是该分支上所有环面 Weil 除子之并。
\item 存在相对余维数 $\ge2$ 的闭子集 $Z\subseteq\mathcal X$，使得 $Z$ 满足以下性质：$Z$ 不包含 $\tilde{\mathcal X}_0$ 的任何环面分层在 $\nu$ 下的像；且对任意点 $x\in\mathcal X\setminus Z$，存在 $x$ 的一个邻域 $\tilde U_x$（在解析拓扑下）、一个 $n+1$ 维仿射环面簇 $Y_x$、$Y_x$ 上由单项式给出的正则函数 $f_x$，以及一个交换图
\[
\begin{tikzcd}
\tilde U_x \arrow[r, "\psi_x"] \arrow[d, "f|_{\tilde U_x}"] & Y_x \arrow[d, "f_x"]\\
D' \arrow[r, "\varphi_x"'] & \C
\end{tikzcd}
\]
其中 $\psi_x$ 与 $\varphi_x$ 是开嵌入，且 $D'\subseteq D$。此外，$f_x$ 在 $Y_x$ 的每个环面除子上都恰好消失到一阶。
\end{enumerate}
\end{definition}

\begin{example}{3.7.2}
取 $\mathcal X$ 由 $\PP^3\times D$ 中的方程 $tf_4+z_0z_1z_2z_3=0$ 定义，其中 $D$ 是以 $t$ 为坐标的圆盘，$f_4$ 是 $\PP^3$ 上的一般齐四次多项式。容易看出 $\mathcal X$ 在轨迹
\[
\{t=f_4=0\}\cap\operatorname{Sing}(\mathcal X_0)
\]
处奇异。由于 $\mathcal X_0$ 是坐标四面体，$\mathcal X_0$ 的奇异轨迹由 $\PP^3$ 的六条坐标直线组成，而 $\mathcal X$ 沿每条这样的直线有四个奇点，共计 24 个奇点。取 $Z=\operatorname{Sing}(\mathcal X)$。于是远离 $Z$ 时投影 $\mathcal X\to D$ 是正规交叉的，这给出环面退化定义中的条件 (4)。容易看出其余条件都满足。
\end{example}

\begin{example}{3.7.3}
设 $\Delta\subseteq M_\R$ 是自反多面体，其对偶为 $\nabla\subseteq N_\R$。如第 4 节那样选取一个整的严格凸分段线性函数 $\check h:M_\R\to\R$。考虑 $\PP_\Delta\times D$ 中由下式定义的族
% [Tutor Note: 式 (3.2) 中首项是 $z^0$（上标 0，即格点 0 对应的截面），不是 $z_0$；指数上的 ȟ 带的是 check（ˇ）而非 tilde；求和只跑遍 $\Delta$ 的格点 $m\in\Delta\cap M$]
\begin{equation}
z^0+t\sum_{m\in\Delta\cap M}a_mt^{\check h'(m)}z^m=0. \tag{3.2}
\end{equation}
这里 $z^m$ 表示由格点 $m\in\Delta\cap M$ 所确定的 $\mathcal O_{\PP_\Delta}(1)$ 的截面，$a_m\in k$ 是一般的系数选取。于是 $z^0$ 是在 $\PP_\Delta$ 的每个环面边界分支上恰好消失一次的截面。

% [原文 p.24]
对 $\check h$ 的大多数选取，这并不定义环面退化：奇点太坏。然而，存在一个自然的环面簇，可以在其中描述这个退化。置
\[
\tilde\Delta:=\{(m,l)\mid m\in\Delta,\ l\ge\check h'(m)\}\subseteq M_\R\oplus\R.
\]
$\check h'$ 为凸的假设蕴涵 $\tilde\Delta$ 是凸的，但当然非紧。设 $\tilde\Sigma$ 是 $\tilde\Delta$ 在 $N_\R\oplus\R$ 中的正规扇。以 $X(\tilde\Sigma)$ 记由扇 $\tilde\Sigma$ 定义的环面簇。于是 $\tilde\Delta$ 是 $X(\tilde\Sigma)$ 上某个线丛 $\mathcal L$ 的牛顿多面体，而形如 $t^{\check h'(m)}z^m$ 的项可解释为该线丛的截面，对应于 $(m,\check h'(m))\in\tilde\Delta\subseteq M_\R\oplus\R$。此外，有自然映射 $X(\tilde\Sigma)\to\mathbb A^1$，由到 $\R$ 的投影定义，而该映射定义了正则函数 $t$。于是式 (3.2) 定义了超曲面 $\mathcal X_\Delta\subseteq X(\tilde\Sigma)$，而 $t$ 定义了映射 $\mathcal X_\Delta\to\mathbb A^1$。这是一个环面退化。
% [Tutor Note: 原文作 refinement $\tilde\Sigma'$ of $\Sigma'$；按上下文，被加细的对象应是 p.24 定义的 $\tilde\Sigma$，而非 $\Sigma'$（$\Sigma'$ 已在 §4/p.12 定义、属另一侧的扇），此处疑为原文笔误；正文照原文排]
不深入细节地说，如第 4 节那样选取 $\nabla$ 的一个星形分解，以及 $\partial\nabla^{\check h}$ 的一个好多面体分解 $\mathcal P$，本质上等同于选取 $\Sigma'$ 的一个具有特别好性质的加细 $\tilde\Sigma'$。这给出部分消解 $X(\tilde\Sigma')\to X(\tilde\Sigma)$。然后 $\mathcal X_\Delta$ 在 $X(\tilde\Sigma')$ 中的真变换 $\mathcal X'_\Delta$ 给出另一个环面退化 $\mathcal X'_\Delta\to\mathbb A^1$。为得到一般纤维是环面簇中超曲面的 MPCP 消解的环面退化，这一步是必要的。证明见 [19]，那里该构造被推广到环面簇中的完全交。\proofend
\end{example}

给定环面退化 $f:\mathcal X\to D$，我们可以如下构造 $f$ 的对偶相交复形 $(B,\mathcal P)$。这里 $B$ 是带奇点的整仿射流形，$\mathcal P$ 是 $B$ 的多面体分解，即把 $B$ 分解为格多面体。事实上，我们将把 $B$ 构造为一些格多面体之并。具体地说，把 $\mathcal X_0$ 的正规化 $\tilde{\mathcal X}_0$ 写成环面簇 $X_i$ 的不交并 $\coprod X_i$，$\nu:\tilde{\mathcal X}_0\to\mathcal X_0$ 为正规化。$\mathcal X_0$ 的分层就是集合
\[
\operatorname{Strata}(\mathcal X_0)=\{\nu(S)\mid S\ \text{是某个}\ X_i\ \text{的环面分层}\}
\]
的元素。这里所谓环面分层，指的是某个 $(\C^*)^n$ 轨道的闭包。
% [Tutor Note: 原文 strata 是 stratum 的复数，指轨道闭包构成的分层，与代数几何里"层（sheaf）"无关；本译本统一译作"分层"]

设 $\{x\}\in\operatorname{Strata}(\mathcal X_0)$ 是零维分层。把定义 3.7.1 的 (4) 应用于 $x$ 的邻域，可知存在环面簇 $Y_x$，使得在 $x$ 的某个邻域内，$f:\mathcal X\to D$ 局部同构于 $f_x:Y_x\to\C$，其中 $f_x$ 由单项式给出。现在，$f_x$ 沿 $Y_x$ 的每个环面除子恰好消失到一阶这一条件，正是 $Y_x$ 为 Gorenstein 的陈述；因此它如习题 3.6.1 那样出现。事实上，取 $M,N$ 如通常，$\operatorname{rank}M=\dim\mathcal X_0$。则存在格多面体 $\sigma_x\subseteq M_\R$，使得 $C(\sigma_x)=\{(rm,r)\mid m\in\sigma,\ r\ge0\}$ 是定义环面簇 $Y_x$ 的锥。
% [Tutor Note: 上一行原文写 $m\in\sigma$，按 Exercise 6.1 中 $C(\sigma)$ 的定义此处应为 $m\in\sigma_x$；已用 550 dpi 裁图复核（p24，该行位于版心中部），确实印作 $\sigma$，属原文笔误，正文照原样排]
正如我们在习题 3.6.1 中看到的，$\mathcal X$ 中 $x$ 的一个小邻域应当向 $B$ 贡献 $\sigma_x$ 的一个拷贝，这为我们的构造提供了动机。现在可以描述如何通过粘合多面体
\[
\{\sigma_x\mid\{x\}\in\operatorname{Strata}(\mathcal X_0)\}
\]
来构造 $B$。我们将在 $\mathcal X_0$ 的每个不可约分支实际上自身正规、从而 $\nu:X_i\to\nu(X_i)$ 是同构的情形下做这件事。读者或许能想象出更一般的构造。在上述正规性假设下，$\sigma_x$ 的面与 $\operatorname{Strata}(\mathcal X_0)$ 中含 $x$ 的元素之间存在一一的、包含关系反转的对应。于是，若 $\sigma_x$ 与 $\sigma_{x'}$ 的面分别对应 $\mathcal X_0$ 的同一个分层，我们就可以把这些面等同起来。需要一些论证来说明这种等同
% [原文 p.25]
可以通过整仿射变换实现，但这同样不难。

作这些等同之后便得到 $B$。于是可以证明

\begin{lemma}{3.7.4}
若 $\mathcal X_0$ 是复 $n$ 维的，则 $B$ 是实 $n$ 维流形。
\end{lemma}

证明见 [22] 中的命题 4.10。
% [Tutor Note: 引理 3.7.4 原文作 "then B is an real n dimensional manifold"，冠词 an 系 a 之笔误；本译本按实 n 维流形译出]

到此为止 $B$ 只是一个拓扑流形，由粘合格多面体而得。置
\[
\mathcal P=\{\sigma\subseteq B\mid \sigma\ \text{是某个零维分层}\ x\ \text{的}\ \sigma_x\ \text{的一个面}\}.
\]
$\mathcal X_0$ 的分层与 $\mathcal P$ 的元素之间存在一一的、包含关系反转的对应。

剩下的只是给 $B$ 一个带奇点的仿射结构。设 $\operatorname{Bar}(\mathcal P)$ 是 $\mathcal P$ 的第一次重心重分，$\Gamma$ 是 $\operatorname{Bar}(\mathcal P)$ 中那些不含 $\mathcal P$ 的顶点、或不与 $\mathcal P$ 的某个极大胞腔的内部相交的单形之并。若置 $B_0:=B\setminus\Gamma$，则可以如下在 $B_0$ 上定义仿射结构。$B_0$ 有开覆盖
\[
\{W_\sigma\mid\sigma\in\mathcal P\ \text{极大}\}\cup\{W_v\mid v\in\mathcal P\ \text{顶点}\},
\]
其中 $W_\sigma=\operatorname{Int}(\sigma)$ 是 $\sigma$ 的内部，而
\[
W_v=\bigcup_{\substack{\tau\in\operatorname{Bar}(\mathcal P)\\ v\in\tau}}\operatorname{Int}(\tau)
\]
是 $v$ 在 $\operatorname{Bar}(\mathcal P)$ 中的（开）星形，正如第 4 节那样。
% [Tutor Note: $\Gamma$ 的定义中"不含 $\mathcal P$ 的顶点"与"不与极大胞腔内部相交"是两条择一的条件（并取所有满足其一的单形），不是同时要求，文字层的 "or" 不可读作 "and"]

现在定义坐标卡。由于 $\mathcal P$ 的极大胞腔形如 $\sigma_x\subseteq M_\R$，此包含给出自然的仿射卡 $\psi_\sigma:\operatorname{Int}(\sigma)\to M_\R$。另一方面，$\mathcal P$ 的顶点 $v$ 对应 $\mathcal X_0$ 的一个余维 0 分层，即某个不可约分支 $X_i$。由于这是紧环面簇，它由 $M_\R$ 中的完全扇 $\Sigma_i$ 定义。此外，$\Sigma_i$ 的 $p$ 维锥与 $\mathcal P$ 中以 $v$ 为顶点的 $p$ 维胞腔之间存在一一对应，因为二者都对应 $\mathcal X_0$ 中容于 $X_i$ 的分层。于是有连续映射
\[
\psi_v:W_v\to M_\R
\]
它把 $W_v\cap\sigma$（对任何以 $v$ 为顶点的 $\sigma\in\mathcal P$）整仿射线性地映入 $\Sigma_i$ 的相应锥中。这样的映射由组合对应关系以及"在每一个胞腔上为整仿射线性"这一要求唯一决定。于是显然这些卡在 $B_0$ 上定义了整仿射结构。这样我们就构造出了 $(B,\mathcal P)$。

\begin{example}{3.7.5}
设 $f:\mathcal X\to D$ 是椭圆曲线退化到 $I_n$ 纤维的退化。则 $B$ 是圆周 $\R/n\Z$，$\mathcal P$ 把它分解为 $n$ 条长度为 1 的线段。
\end{example}

\begin{example}{3.7.6}
承接例 3.7.3，从环面退化 $\mathcal X'_\Delta\to\mathbb A^1$ 构造出的对偶相交复形，就是第 4 节中在 $B=\partial\nabla^{\check h}$ 上构造的带奇点仿射流形结构与多面体分解 $(B,\mathcal P)$。这并不特别难证明。同样，证明以及更一般的完全交情形见 [19]。
\end{example}
